\documentclass[10pt,a4paper]{amsart}
\usepackage{amsmath,amssymb,amsthm,graphicx}
% For commutative diagrams we use Xy-pic
\usepackage[all]{xy}

% We use 2cell for 2-commutative diagrams.
\xyoption{2cell}
% Theorem environments.
%
\theoremstyle{plain}
\newtheorem{theorem}[subsection]{Theorem}
\newtheorem{proposition}[subsection]{Proposition}
\newtheorem{lemma}[subsection]{Lemma}

\theoremstyle{definition}
\newtheorem{definition}[subsection]{Definition}
\newtheorem{example}[subsection]{Example}
\newtheorem{exercise}[subsection]{Exercise}
\newtheorem{situation}[subsection]{Situation}

\theoremstyle{remark}
\newtheorem{remark}[subsection]{Remark}
\newtheorem{remarks}[subsection]{Remarks}

\numberwithin{equation}{subsection}

% Macros
%
\def\lim{\mathop{\mathrm{lim}}\nolimits}
\def\colim{\mathop{\mathrm{colim}}\nolimits}
\def\Spec{\mathop{\mathrm{Spec}}}
\def\Hom{\mathop{\mathrm{Hom}}\nolimits}
\def\Ext{\mathop{\mathrm{Ext}}\nolimits}
\def\SheafHom{\mathop{\mathcal{H}\!\mathit{om}}\nolimits}
\def\SheafExt{\mathop{\mathcal{E}\!\mathit{xt}}\nolimits}
\def\Sch{\mathit{Sch}}
\def\Mor{\mathop{\mathrm{Mor}}\nolimits}
\def\Ob{\mathop{\mathrm{Ob}}\nolimits}
\def\Sh{\mathop{\mathit{Sh}}\nolimits}
\def\NL{\mathop{N\!L}\nolimits}
\def\CH{\mathop{\mathrm{CH}}\nolimits}
\def\proetale{{pro\text{-}\acute{e}tale}}
\def\etale{{\acute{e}tale}}
\def\QCoh{\mathit{QCoh}}
\def\Ker{\mathop{\mathrm{Ker}}}
\def\Im{\mathop{\mathrm{Im}}}
\def\Coker{\mathop{\mathrm{Coker}}}
\def\Coim{\mathop{\mathrm{Coim}}}

% Boxtimes
%
\DeclareMathSymbol{\boxtimes}{\mathbin}{AMSa}{"02}

%
% Macros for moduli stacks/spaces
%
\def\QCohstack{\mathcal{QC}\!\mathit{oh}}
\def\Cohstack{\mathcal{C}\!\mathit{oh}}
\def\Spacesstack{\mathcal{S}\!\mathit{paces}}
\def\Quotfunctor{\mathrm{Quot}}
\def\Hilbfunctor{\mathrm{Hilb}}
\def\Curvesstack{\mathcal{C}\!\mathit{urves}}
\def\Polarizedstack{\mathcal{P}\!\mathit{olarized}}
\def\Complexesstack{\mathcal{C}\!\mathit{omplexes}}
% \Pic is the operator that assigns to X its picard group, usage \Pic(X)
% \Picardstack_{X/B} denotes the Picard stack of X over B
% \Picardfunctor_{X/B} denotes the Picard functor of X over B
\def\Pic{\mathop{\mathrm{Pic}}\nolimits}
\def\Picardstack{\mathcal{P}\!\mathit{ic}}
\def\Picardfunctor{\mathrm{Pic}}
\def\Deformationcategory{\mathcal{D}\!\mathit{ef}}
\usepackage{cleveref}
\crefname{originaldefinition}{original definition}{original definitions}
\Crefname{originaldefinition}{Original definition}{Original definitions}
\makeatletter
\newlabel{definition-cech-complex}{{03OL}{source}{}{sourcepointer-1}{}}
\newlabel{definition-cech-complex@cref}{{[originaldefinition][0][]03OL}{[1][0][]source}{}{}{}}
\makeatother
\crefname{originallemma}{original lemma}{original lemmas}
\Crefname{originallemma}{Original lemma}{Original lemmas}
\makeatletter
\newlabel{lemma-hom-injective}{{03OR}{source}{}{sourcepointer-2}{}}
\newlabel{lemma-hom-injective@cref}{{[originallemma][0][]03OR}{[1][0][]source}{}{}{}}
\makeatother
\crefname{originallemma}{original lemma}{original lemmas}
\Crefname{originallemma}{Original lemma}{Original lemmas}
\makeatletter
\newlabel{sites-cohomology-lemma-cech-cohomology-derived-presheaves}{{03AU}{source}{}{sourcepointer-3}{}}
\newlabel{sites-cohomology-lemma-cech-cohomology-derived-presheaves@cref}{{[originallemma][0][]03AU}{[1][0][]source}{}{}{}}
\makeatother
\title{Cech cohomology of presheaves computes the right derived functors of H\textasciicircum{}0-check}
\author{The Stacks Project Authors}
\date{stacks-project:etale-cohomology.tex:local-snapshot-20261011; GNU FDL 1.2 or later}
\begin{document}
\maketitle
\noindent\textit{Editorial note.} All delivered mathematical text is the author's own, copied from the pinned Stacks Project chapter file etale-cohomology.tex; no statement or proof step was rewritten, reordered or removed.\par\medskip
\noindent\textit{Editorial note.} The preamble is the chapter's own native preamble excerpt used for this source: the Xy-pic setup at lines 14-18 and the theorem declarations with the mathematical macros at lines 154-217 of sources/stacks-etale-cohomology/preamble.tex. The quoted class line and the xr-hyper externaldocument block are not included because the assembler supplies a plain amsart wrapper; no macro was invented or redefined.\par\medskip
\noindent\textit{Editorial note.} The same-chapter references are located at their \textbackslash{}label line inside the pinned chapter file; the cross-chapter pointer is located in the project's own tag list staged in the snapshot (sources/stacks-chow/tags.txt), at the verbatim line carrying its canonical tag.\par\medskip
\noindent\textit{Allowed standard prerequisites (editorial).}\begin{itemize}
\item Cech complex and Cech cohomology of a family in a presheaf of abelian groups. Same chapter, Definition environment labelled definition-cech-complex (source line 1477, Stacks tag 03OL); the statement opens by referring to its notation and assumptions. Supplies the functors H\textasciicircum{}p-check(U, -) on PAb(C); the comparison with derived functors is the counted content.
\item Universal delta-functors and effaceability. Standard homological algebra; the proof uses that an effaceable delta-functor is universal and that two universal delta-functors agreeing in degree 0 agree everywhere. Supplies the two general facts combined in the proof; the identification of effaceability with the previous lemma is the counted step.
\end{itemize}\medskip
\ifdefined\hypertarget\hypertarget{sourcepointer-1}{}\fi
\noindent\textit{Source reference.} Original definition 03OL: sources/stacks-etale-cohomology/pinned-original-source.tex, source lines 1477--1477. Same-chapter pointer named in the first line of the counted statement (line 1725) as the source of notation and assumptions; the comparison of delta-functors is retained verbatim in the counted span.\par\medskip
\ifdefined\hypertarget\hypertarget{sourcepointer-2}{}\fi
\noindent\textit{Source reference.} Original lemma 03OR: sources/stacks-etale-cohomology/pinned-original-source.tex, source lines 1702--1702. Same-chapter pointer invoked at proof line 1731 for effaceability of the Cech functors; it supplies that input only, while the universality and degree-zero comparison are retained in the counted span.\par\medskip
\ifdefined\hypertarget\hypertarget{sourcepointer-3}{}\fi
\noindent\textit{Source reference.} Original lemma 03AU: sources/stacks-chow/tags.txt, source lines 4323--4323. Cross-chapter 'for more details' pointer at proof lines 1736-1738, invoked after the argument is complete; the retained proof is self-contained in the two delta-functor facts.\par\medskip
\begin{theorem}
\label{theorem-cech-derived}
Notation and assumptions as in Definition \ref{definition-cech-complex}.
On $\textit{PAb}(\mathcal{C})$ the functors $\check{H}^p(\mathcal{U}, -)$ are
the right derived functors of $\check{H}^0(\mathcal{U}, -)$.
\end{theorem}

\begin{proof}
By the Lemma \ref{lemma-hom-injective}, the functors
$\check H^p(\mathcal{U}, -)$ are universal
$\delta$-functors since they are effaceable.
So are the right derived functors of $\check H^0(\mathcal{U}, -)$. Since they
agree in degree $0$, they agree by the universal property of universal
$\delta$-functors. For more details see
Cohomology on Sites,
Lemma \ref{sites-cohomology-lemma-cech-cohomology-derived-presheaves}.
\end{proof}

\end{document}
